\section{Vlastné hodnoty, vlastné vektory}

Nech $A \colon V \to V$ je lineárna transformácia konečnorozmerného priestoru, $\dim(V) = n$.

Hľadáme bázu $X = (\ora{x}_1, \dots, \ora{x}_n)$ takú, že $[A]_{XX}$ je diagonálna matica.

Čo to znamená?
\[
  [A]_{XX}
  = \begin{pmatrix}
      [A\ora{x}_1]_X & \cdots & [A\ora{x}_n]_X
    \end{pmatrix}
\]
Toto má byť diagonálna matica, teda
\[
  [A\ora{x}_1]_X
  = \begin{pmatrix} \lambda_1 \\ 0 \\ \vdots \\ 0 \end{pmatrix}, \quad
  [A\ora{x}_2]_X
  = \begin{pmatrix} 0 \\ \lambda_2 \\ \vdots \\ 0 \end{pmatrix}, \quad \dots \quad
  [A\ora{x}_n]_X
  = \begin{pmatrix} 0 \\ 0 \\ \vdots \\ \lambda_n \end{pmatrix}
\]
(toto sú súradnice v~báze $X$). Teda
\[
  [A\ora{x}_i]_X
  = \begin{pmatrix} 0 \\ \vdots \\ 0 \\ \lambda_i \\ 0 \\ \vdots \\ 0 \end{pmatrix}
  \leftarrow \text{v~$i$-tom riadku}
\]
a teda
\[
  A\ora{x}_i
  = 0 \cdot \ora{x}_1 + \dots + 0 + \lambda_i \ora{x}_i + 0 + \dots + 0 \cdot \ora{x}_n
  = \lambda_i \ora{x}_i.
\]

Teda každý vektor $\ora{x}$, ktorý sa vyskytuje v~tej báze $X$, má vlastnosť
\[
  A\ora{x} = \lambda\ora{x}
  \quad \text{pre nejaký skalár } \lambda \in K.
\]
Zároveň $\ora{x} \neq \ora{o}$ (lebo $\ora{o}$ nie je v~žiadnej báze).

\begin{definicia}
\label{def:vlastna-hodnota}
Nech $A \colon V \to V$ je lineárna transformácia. Ak
\[
  A\ora{x} = \lambda\ora{x}, \quad \ora{x} \neq \ora{o}, \quad \lambda \in K,
\]
potom hovoríme, že $\lambda$ je \emph{vlastná hodnota} $A$
a~$\ora{x}$ je k~nej prislúchajúci \emph{vlastný vektor}.
\end{definicia}

Podľa horeuvedených úvah teda dostávame:

\begin{tvrdenie}
\label{tvrden:diag-vlastne-vektory}
Nech $A \colon V \to V$ je lineárna transformácia konečnorozmerného vektorového priestoru $V$, $X$ je báza $V$.
Potom $[A]_{XX}$ je diagonálna práve vtedy, keď $X$ pozostáva z~vlastných vektorov $A$.
\end{tvrdenie}

\begin{priklad}
Uvažujme lineárnu transformáciu „osová súmernosť podľa fixnej priamky $P$ v~rovine $S$ s~počiatkom''.

% TikZ diagram for the reflection example
\begin{center}
\begin{tikzpicture}[scale=1.2]
  % Line P
  \draw[blue, thick] (-1,-1.5) -- (2,2.5) node[above right]{$P$};
  % Original vector
  \draw[->, thick] (0,0) -- (0.5,1.5) node[above left]{$\ora{x}$};
  % Reflected vector
  \draw[->, thick, green!60!black] (0,0) -- (1.5,0.5) node[right]{$O_P\ora{x}$};
\end{tikzpicture}
\end{center}

Aké sú vlastné vektory tohto lineárneho zobrazenia $O_P$?

Sú dvoch druhov:
\begin{enumerate}
  \item Tie, ktoré ležia na priamke $P$; platí pre ne
    $O_P(\ora{x}) = \ora{x} = 1 \cdot \ora{x}$,
    teda prislúchajú ku vlastnej hodnote $1$.

  \item Tie, ktoré sú \underline{kolmé} na $P$; platí pre ne
    $O_P(\ora{x}) = (-1)\ora{x}$,
    teda prislúchajú vlastnej hodnote $(-1)$.
\end{enumerate}

Iné vlastné hodnoty ani vektory $O_P$ zrejme nemá.
\end{priklad}

Rozmýšľajme teraz o~nejakej lineárnej transformácii $A \colon V \to V$.
Ako nájsť vlastné hodnoty a~vlastné vektory $A$?

Predpokladajme najprv, že o~nejakom $\lambda \in K$ vieme, že je vlastnou hodnotou $A$, to znamená
\[
  A\ora{x} = \lambda\ora{x}, \quad \ora{x} \neq \ora{o}.
\]
Navzájom ekvivalentné výroky o~$\ora{x}$:
\begin{align*}
  A\ora{x} - \lambda\ora{x} &= \ora{o} \\
  A\ora{x} - \lambda I\,\ora{x} &= \ora{o} \\
  (A - \lambda I)\ora{x} &= \ora{o} \\
  \ora{x} &\in \Ker(A - \lambda I)
\end{align*}

Ak teda vieme, že $\lambda$ je vlastná hodnota, vlastné vektory prislúchajúce k~$\lambda$ môžeme určiť ako prvky jadra akejsi lineárnej transformácie.

\begin{priklad}
Uvažujme lineárnu transformáciu vektorového priestoru $\R^2$ danú maticou
\[
  A = \begin{pmatrix} 1 & 1 \\ 3 & -1 \end{pmatrix}
\]
v~kanonickej báze $\R^2$.

Prezradím, že vlastné hodnoty tejto lineárnej transformácie sú
\[
  \lambda_1 = 2, \quad \lambda_2 = -2.
\]

Nájdime vlastný vektor prislúchajúci k~$\lambda_1 = 2$.
\[
  A - \lambda_1 I
  = \begin{pmatrix} 1 & 1 \\ 3 & -1 \end{pmatrix}
  - \begin{pmatrix} 2 & 0 \\ 0 & 2 \end{pmatrix}
  = \begin{pmatrix} -1 & 1 \\ 3 & -3 \end{pmatrix}
\]

Hľadáme $\Ker(A - \lambda_1 I)$:
\[
  \left(\begin{array}{cc|c}
    -1 & 1 & 0 \\
     3 & -3 & 0
  \end{array}\right)
  \sim
  \left(\begin{array}{cc|c}
    -1 & 1 & 0 \\
     0 & 0 & 0
  \end{array}\right)
\]

Položme $x_2 = s$. Potom $-x_1 + x_2 = 0$, teda $-x_1 + s = 0$, čiže $x_1 = s$. Dostávame
\[
  \Ker(A - \lambda_1 I)
  = \{(s, s) \mid s \in \R\}.
\]

Napríklad $(1,1)$ je vlastný vektor prislúchajúci k~$2$. Overenie:
\[
  \begin{pmatrix} 1 & 1 \\ 3 & -1 \end{pmatrix}
  \begin{pmatrix} 1 \\ 1 \end{pmatrix}
  = \begin{pmatrix} 2 \\ 2 \end{pmatrix}
  = 2 \cdot \begin{pmatrix} 1 \\ 1 \end{pmatrix}
  = \lambda_1 \cdot \begin{pmatrix} 1 \\ 1 \end{pmatrix}.
\]

Keď urobíme to isté pre tú druhú vlastnú hodnotu, dostaneme
\[
  A - \lambda_2 I
  = \begin{pmatrix} 3 & 1 \\ 3 & 1 \end{pmatrix}.
\]

Hľadáme $\Ker(A - \lambda_2 I)$:
\[
  \left(\begin{array}{cc|c}
    3 & 1 & 0 \\
    3 & 1 & 0
  \end{array}\right)
  \sim
  \left(\begin{array}{cc|c}
    3 & 1 & 0 \\
    0 & 0 & 0
  \end{array}\right)
\]

Dostávame $3x_1 + x_2 = 0$. Položme $x_2 = t$, potom $3x_1 + t = 0$, čiže $x_1 = -\frac{1}{3}t$. Teda
\[
  \Ker(A - \lambda_2 I)
  = \left\{\left(-\tfrac{1}{3}t,\; t\right) \;\middle|\; t \in \R\right\}.
\]

Čiže jeden z~vlastných vektorov prislúchajúcich k~$\lambda_2 = -2$ je (pre $t = 1$) vektor $\left(-\frac{1}{3},\; 1\right)$.
\end{priklad}
